% GENERATED FILE. Edit canonical lessons, then run npm run generate:books.
% canonical-source-hash: fnv1a64:8f361d017cd65112
% canonical-lessons: AR-C272-mudhish, AR-C272-mutaakkid, AR-C272-muthir, AR-C272-muaqqad, AR-C272-maqul

\chapter{Amazing, Sure, Exciting, Complicated, Reasonable}
\label{ch:ar-amazing-sure-exciting-complicated-reasonable}
\hlchaptermodality{272}

\begin{chapteropening}
\textbf{By the end of this chapter:} \emph{I can say amazing, sure, exciting, complicated and reasonable.}

Complete the last lesson of chapter 272: I can say amazing, sure, exciting, complicated and reasonable.
\end{chapteropening}

\section[\texorpdfstring{\ar{مدهش} (mudhish)}{مدهش (mudhish)}]{\ar{مدهش} (mudhish) — amazing}

\label{lesson:AR-C272-mudhish}



\begin{quote}
Before the new one: say the Arabic for pink, then the Arabic for purple.
\end{quote}

\subsection*{You'll want to know: \ar{مدهش}}
\textbf{\ar{مدهش}} — \emph{mudhish} — \textquotedblleft{}amazing\textquotedblright{}.

\begin{cousinweb}
\textbf{\ar{مدهش}} (\emph{mudhish}) is built on the native Arabic root \emph{d-h-sh}, \textquotedblleft{}to be astonished\textquotedblright{} (Lane). \emph{Mudhish} is the Form IV participle, \textquotedblleft{}astonishing\textquotedblright{}.
\end{cousinweb}

\begin{grammarlens}[title={What you've built}]
One more word for this chapter.
\end{grammarlens}

\subsection*{Your turn}
\begin{itemize}
\raggedright
  \item \emph{Say it:} \emph{mudhish}
  \item \emph{Say it:} \emph{mudhish}, once more
  \item \emph{Say it:} say \emph{mudhish}
  \item \emph{Recall:} say the Arabic for pink, then the Arabic for purple, then say \emph{mudhish} again
\end{itemize}

\begin{tcolorbox}[breakable,colback=teal!4,colframe=teal!35!black,title={Before you move on}]
What does \ar{مدهش} mean? (\textquotedblleft{}Amazing\textquotedblright{}.) Say it once more.
\end{tcolorbox}
\section[\texorpdfstring{\ar{متأكد} (mutaʾakkid)}{متأكد (mutaʾakkid)}]{\ar{متأكد} (mutaʾakkid) — sure, certain}

\label{lesson:AR-C272-mutaakkid}



\begin{quote}
Before the new one: say the Arabic for purple, then the Arabic for amazing.
\end{quote}

\subsection*{You'll want to know: \ar{متأكد}}
\textbf{\ar{متأكد}} — \emph{mutaʾakkid} — \textquotedblleft{}sure, certain\textquotedblright{}.

\begin{cousinweb}
\textbf{\ar{متأكد}} (\emph{mutaʾakkid}) is built on the native Arabic root \emph{ʾ-k-d}, \textquotedblleft{}to make firm\textquotedblright{} (Lane). \emph{Mutaʾakkid} is the Form V participle, \textquotedblleft{}made firm\textquotedblright{}, hence \textquotedblleft{}sure\textquotedblright{}.
\end{cousinweb}

\begin{grammarlens}[title={What you've built}]
One more word for this chapter.
\end{grammarlens}

\subsection*{Your turn}
\begin{itemize}
\raggedright
  \item \emph{Say it:} \emph{mutaʾakkid}
  \item \emph{Say it:} \emph{mutaʾakkid}, once more
  \item \emph{Say it:} say \emph{mutaʾakkid}
  \item \emph{Recall:} say the Arabic for purple, then the Arabic for amazing, then say \emph{mutaʾakkid} again
\end{itemize}

\begin{tcolorbox}[breakable,colback=teal!4,colframe=teal!35!black,title={Before you move on}]
What does \ar{متأكد} mean? (\textquotedblleft{}Sure, certain\textquotedblright{}.) Say it once more.
\end{tcolorbox}
\section[\texorpdfstring{\ar{مثير} (muthīr)}{مثير (muthīr)}]{\ar{مثير} (muthīr) — exciting}

\label{lesson:AR-C272-muthir}



\begin{quote}
Before the new one: say the Arabic for amazing, then the Arabic for sure.
\end{quote}

\subsection*{You'll want to know: \ar{مثير}}
\textbf{\ar{مثير}} — \emph{muthīr} — \textquotedblleft{}exciting\textquotedblright{}.

\begin{cousinweb}
\textbf{\ar{مثير}} (\emph{muthīr}) is built on the native Arabic root \emph{th-w-r}, \textquotedblleft{}to rise up; a bull\textquotedblright{} (Lane). \emph{Muthīr} is the Form IV participle, \textquotedblleft{}stirring up\textquotedblright{}.
\end{cousinweb}

\begin{grammarlens}[title={What you've built}]
One more word for this chapter.
\end{grammarlens}

\subsection*{Your turn}
\begin{itemize}
\raggedright
  \item \emph{Say it:} \emph{muthīr}
  \item \emph{Say it:} \emph{muthīr}, once more
  \item \emph{Say it:} say \emph{muthīr}
  \item \emph{Recall:} say the Arabic for amazing, then the Arabic for sure, then say \emph{muthīr} again
\end{itemize}

\begin{tcolorbox}[breakable,colback=teal!4,colframe=teal!35!black,title={Before you move on}]
What does \ar{مثير} mean? (\textquotedblleft{}Exciting\textquotedblright{}.) Say it once more.
\end{tcolorbox}
\section[\texorpdfstring{\ar{معقد} (muʿaqqad)}{معقد (muʿaqqad)}]{\ar{معقد} (muʿaqqad) — complicated}

\label{lesson:AR-C272-muaqqad}



\begin{quote}
Before the new one: say the Arabic for sure, then the Arabic for exciting.
\end{quote}

\subsection*{You'll want to know: \ar{معقد}}
\textbf{\ar{معقد}} — \emph{muʿaqqad} — \textquotedblleft{}complicated\textquotedblright{}.

\begin{cousinweb}
\textbf{\ar{معقد}} (\emph{muʿaqqad}) is built on the native Arabic root \emph{ʿ-q-d}, \textquotedblleft{}to tie a knot\textquotedblright{} (Lane). \emph{Muʿaqqad} is the Form II passive participle, \textquotedblleft{}knotted\textquotedblright{}.
\end{cousinweb}

\begin{grammarlens}[title={What you've built}]
One more word for this chapter.
\end{grammarlens}

\subsection*{Your turn}
\begin{itemize}
\raggedright
  \item \emph{Say it:} \emph{muʿaqqad}
  \item \emph{Say it:} \emph{muʿaqqad}, once more
  \item \emph{Say it:} say \emph{muʿaqqad}
  \item \emph{Recall:} say the Arabic for sure, then the Arabic for exciting, then say \emph{muʿaqqad} again
\end{itemize}

\begin{tcolorbox}[breakable,colback=teal!4,colframe=teal!35!black,title={Before you move on}]
What does \ar{معقد} mean? (\textquotedblleft{}Complicated\textquotedblright{}.) Say it once more.
\end{tcolorbox}
\section[\texorpdfstring{\ar{معقول} (maʿqūl)}{معقول (maʿqūl)}]{\ar{معقول} (maʿqūl) — reasonable}

\label{lesson:AR-C272-maqul}



\begin{quote}
Before the new one: say the Arabic for exciting, then the Arabic for complicated.
\end{quote}

\subsection*{You'll want to know: \ar{معقول}}
\textbf{\ar{معقول}} — \emph{maʿqūl} — \textquotedblleft{}reasonable\textquotedblright{}.

\begin{cousinweb}
\textbf{\ar{معقول}} (\emph{maʿqūl}) is built on the native Arabic root \emph{ʿ-q-l}, \textquotedblleft{}to tie; to understand\textquotedblright{} (Lane). \emph{Maʿqūl} is the passive participle, \textquotedblleft{}understood\textquotedblright{}, hence \textquotedblleft{}reasonable\textquotedblright{}.
\end{cousinweb}

\begin{grammarlens}[title={What you've built}]
One more word for this chapter.
\end{grammarlens}

\subsection*{Your turn}
\begin{itemize}
\raggedright
  \item \emph{Say it:} \emph{maʿqūl}
  \item \emph{Say it:} \emph{maʿqūl}, once more
  \item \emph{Say it:} say \emph{maʿqūl}
  \item \emph{Recall:} say the Arabic for exciting, then the Arabic for complicated, then say \emph{maʿqūl} again
\end{itemize}

\begin{tcolorbox}[breakable,colback=teal!4,colframe=teal!35!black,title={Before you move on}]
What does \ar{معقول} mean? (\textquotedblleft{}Reasonable\textquotedblright{}.) Say it once more.
\end{tcolorbox}
